\section{模型假设}

\subsection{问题一的模型假设}

\begin{enumerate}
  \item 问题一题干原文要求“计算多边形定位区域直径”，故本问将每次检测所得示向度的误差区域，视为从检测点出发、向前无限延伸的两条射线所夹的扇形区域，定位区域只取这些扇形区域的公共部分，不再使用有效接收半径或半径为 $1800\ \mathrm{m}$ 的目标圆形区域对其进行裁剪；
  \item 主要讨论交会有效的情形，即多个角度误差区域的公共部分非空、有界且面积不为零。若公共部分为空或无界，则不存在题目所述的多边形直径。
\end{enumerate}

\subsection{问题二的模型假设}

\begin{enumerate}
  \item 示向度误差按题设的有界误差处理，不预设其服从特定概率分布。连续选点中的等方差局部扰动仅用于构造 Fisher 信息替代指标，不改变最终的定位不确定性半径评价标准；
  \item 为保持源位置区域的凸性，不从首次检测所得示向度对应的区域中扣除半径为 $5\ \mathrm{m}$ 的“近距离”圆孔。该处理只可能放大可行域，不会排除真实源；
  \item 对于无法保证接收的第二测点，如果在该点未收到信号且定位区域也没有有效缩小，就视为不利情况，暂时不利用无信号结果可能提供的排除信息。
\end{enumerate}

\subsection{问题三的模型假设}

\begin{enumerate}
  \item 对每个频道分别维护源位置可行域；已经获得示向度的频道在后续扫描中不再承担“发现新源”的作用，其重复检测仅用于缩小该频道已有源的定位区域；
  \item 默认策略不利用单次无信号结果切除源位置区域。一个频道只有在八个保证覆盖驻留点均完成检测、且从未返回示向度或近距离结果时，才被归入无源频道集合；
  \item 补充定位阶段采用有限源位置场景、有限误差端点和有限次数的数值选点。该模型用于产生可执行的较优测点，不假定其达到连续空间中的全局最优；
  \item 多源调度采用“先形成清除队列、再优化队列中心访问顺序”的分阶段近似。旅行商模型优化的是各中心之间的开放路径，不把定位过程中已经发生的移动及从当前位置到首个清除点的距离纳入同一全局目标；
  \item 对尚未完成可靠定位的源使用有限探针集合近似连续清除区域。仅在完整覆盖该源最小包围圆所需的探针数不超过程序上限时，使用其确定性覆盖结论；其余情形作为有限试探策略处理。
\end{enumerate}

\subsection{问题四的模型假设}

\begin{enumerate}
  \item 有效示向度继续用于更新连续的源位置外近似；单次无信号结果只用于筛除与其不相容的有限可能状态，不直接从凸多边形定位区域中挖去非凸部分；
  \item 为控制在线计算量，源位置从当前定位区域中选取不超过 $12$ 个代表点，接收半径取 $1000$、$1250$、$1500\ \mathrm{m}$，定向角以 $10^\circ$ 间隔离散，并补充各历史检测点对应的可见边界方向；该有限化只用于比较下一步行动，不作为连续空间的完整刻画；
  \item 在缺少全向源和定向源比例信息时，两类源的初始总相对权重均取 $0.5$，再在各自的有限可能状态中均分；这些权重仅服务于行动时间的相对评价，不解释为统计调查所得概率；
  \item 获得新示向度后，清除覆盖点数近似随定位区域最小包围圆半径的平方变化；该面积比例用于预估后续行动时间，实际是否清除以及定位区域更新仍以模拟器真实返回结果为准；
  \item 三角网平移量以及最近邻加 $2$-opt 的扫描顺序是满足覆盖要求的有限数值选择，只主张当前点集具有定向发现保证，不假定驻留点数量或扫描路径达到全局最优。
\end{enumerate}
